\documentclass[conference]{IEEEtran}
\IEEEoverridecommandlockouts
% The preceding line is only needed to identify funding in the first footnote. If that is unneeded, please comment it out.
%Template version as of 6/27/2024

\usepackage{cite}
\usepackage{amsmath,amssymb,amsfonts}
\usepackage{algorithmic}
\usepackage{graphicx}
\usepackage{textcomp}
\usepackage{xcolor}
\def\BibTeX{{\rm B\kern-.05em{\sc i\kern-.025em b}\kern-.08em
    T\kern-.1667em\lower.7ex\hbox{E}\kern-.125emX}}
    
\begin{document}

\title{Thai Poetry Rhyme Verification using Hybrid Rules Based Approach for Education}

\author{\IEEEauthorblockN{1\textsuperscript{st} Warit Yuvaniyama}
\IEEEauthorblockA{
\textit{School of Information, Computer,} \\
\textit{and Communication Technology} \\
\textit{Sirindhorn International Institute} \\
\textit{of Technology, Thammasat University}\\
Pathum Thani 12120, Thailand \\
https://orcid.org/0009-0007-1630-4286}
\and
\IEEEauthorblockN{2\textsuperscript{nd} Athit Phinyachok}
\IEEEauthorblockA{\textit{School of Information, Computer,} \\
\textit{and Communication Technology} \\
\textit{Sirindhorn International Institute} \\
\textit{of Technology, Thammasat University}\\
Pathum Thani 12120, Thailand \\
https://orcid.org/0009-0000-3783-8104}
\and
\IEEEauthorblockN{3\textsuperscript{rd} Gunn Sanguankittipun}
\IEEEauthorblockA{\textit{School of Information, Computer,} \\
\textit{and Communication Technology} \\
\textit{Sirindhorn International Institute} \\
\textit{of Technology, Thammasat University}\\
Pathum Thani 12120, Thailand \\
https://orcid.org/0009-0005-8812-4992}
\and
\IEEEauthorblockN{4\textsuperscript{th} Leo Colombo}
\IEEEauthorblockA{\textit{School of Information, Computer,} \\
\textit{and Communication Technology} \\
\textit{Sirindhorn International Institute} \\
\textit{of Technology, Thammasat University}\\
Pathum Thani 12120, Thailand \\
https://orcid.org/0009-0004-9133-5398}
\and
\IEEEauthorblockN{5\textsuperscript{th} Prachya Boonkwan*}
\IEEEauthorblockA{\textit{School of Information, Computer,} \\
\textit{and Communication Technology} \\
\textit{Sirindhorn International Institute} \\
\textit{of Technology, Thammasat University}\\
Pathum Thani 12120, Thailand \\
https://orcid.org/0000-0003-1405-1071*}
}

\maketitle

\begin{abstract}
% [Drafting note: Keep this under 250 words. Write this LAST.]
% 1. Context/Problem: LLMs struggle to generate Thai poetry (Klon-8) due to strict phonetic, rhyme, and meter constraints. Educators also lack reliable automated tools.
% 2. Gap: Existing rule-based and G2P checkers suffer from low recall, false positives on vowel lengths, and blind spots with complex Thai morphology.
% 3. Solution: We propose an improved hybrid rule-based verification algorithm integrated into PyThaiNLP 5.3.5, alongside a G2P-override variant (Klonpad).
% 4. Evaluation: Evaluated on a massive gold-standard dataset of 36,475 classical stanzas (145,709 rhyme checks) with targeted data augmentation for edge cases.
% 5. Impact: The model achieves 88.5\% gold recall and near 100\% precision on curated negatives, serving as both an educational web tool and a deterministic reward model for RLHF in LLM training.
\end{abstract}

\begin{IEEEkeywords}
Thai NLP, Poetry Generation, Rule-based Verification, Natural Language Processing, Reinforcement Learning, Educational Technology
\end{IEEEkeywords}

\section{Introduction}
% [Drafting note: Set the stage and tell the story of WHY this matters.]
\begin{itemize}
    \item \textbf{Context:} Introduce Thai Klon-Paed (กลอนแปด) and its strict structural rules (Sumpus, Mattra, Sara).
    \item \textbf{The Core Problem:} Generative AI and LLMs are probabilistic, making them inherently bad at strict deterministic constraints like syllables and rhymes.
    \item \textbf{The Dual Need:} 
        \begin{itemize}
            \item \textit{Research:} To train an LLM to write poetry (e.g., via RLHF/RLAIF), we need a flawless, programmatic reward environment to score rhymes.
            \item \textit{Education:} Teachers need zero-hallucination, interactive grading tools for students.
        \end{itemize}
    \item \textbf{Limitations of Prior Art:} Previous checkers (PyThaiNLP 5.0.1) and G2P models (Kongfha) fail on edge cases, drop stanzas, or are computationally too slow for RL loops.
    \item \textbf{Contributions:} 
        \begin{itemize}
            \item Improved core verification algorithms (\texttt{check\_sara}, \texttt{check\_marttra}, \texttt{is\_sumpus}).
            \item Creation of a 36,475-stanza gold-standard evaluation dataset.
            \item Integration into PyThaiNLP 5.3.5 and an open-source Hugging Face educational application.
        \end{itemize}
\end{itemize}

\section{Background and Related Work}
% [Drafting note: Briefly explain the rules and the baseline models.]
\subsection{Linguistic Rules of Klon-Paed}
\begin{itemize}
    \item Briefly define the requirements for a valid rhyme in Thai: matching vowel sound (\texttt{check\_sara}) and matching final consonant class (\texttt{check\_marttra}).
    \item Define inter-stanza (rX) and intra-stanza (r1, r2, r3) connections.
\end{itemize}

\subsection{Existing Verification Systems}
\begin{itemize}
    \item \textbf{Model A (Old PyThaiNLP 5.0.1):} Dictionary-based, lacks the r3 rule, fails on complex morphology.
    \item \textbf{Model C (Kongfha / tltk G2P):} Romanization-based. Extremely slow (single-threaded) and collapses short/long vowels, leading to false positives.
\end{itemize}

\section{Methodology: Hybrid Verification Architecture}
% [Drafting note: Explain YOUR algorithm and the two variants you built.]
\subsection{Core Rule-Based Improvements (Model B)}
\begin{itemize}
    \item Explain the mechanics of Model B (now PyThaiNLP 5.3.5).
    \item \textbf{Phonetic Normalization:} Handling transformed/reduced vowels (สระประสม, สระลดรูป).
    \item \textbf{Edge Case Handling:} 
        \begin{itemize}
            \item \texttt{handle\_karun\_sound\_silence}: Stripping silenced characters (e.g., "โอห์ม" $\rightarrow$ "โอ").
            \item \texttt{\_is\_true\_final}: Differentiating genuine finals (-ย, -ล) from clusters or digraphs.
        \end{itemize}
\end{itemize}

\subsection{Dictionary-Enhanced Verification (Model D)}
\begin{itemize}
    \item Explain the Klonpad variant. While Model B acts as a general-purpose oracle, it has "blind spots" (e.g., the silent 'ร' in "เพชร").
    \item Introduce the G2P Override Dictionary (built on Phra Aphai Mani) to fix these specific phonetic misreadings.
    \item \textbf{Ablation Study of Fallbacks:} Briefly discuss why D\_ssg (SSG fallback) outperformed D\_w2p (W2P fallback). W2P introduces hallucinations on short words, hurting overall accuracy.
\end{itemize}

\section{Dataset Construction and Augmentation}
% [Drafting note: Showcase the rigorous testing environment.]
\subsection{Gold-Standard Corpus Compilation}
\begin{itemize}
    \item Data sourced from vajirayana.org (Khun Chang Khun Phaen, Phra Aphai Mani, etc.).
    \item Cleaning pipeline resulting in 191 files, 36,475 stanzas, and 145,709 total rhyme checks.
    \item Highlight the data-completeness fix (retaining $\geq 10$ syllable waks) that restored inter-stanza (rX) chains.
\end{itemize}

\subsection{Targeted Data Augmentation}
\begin{itemize}
    \item Explain the need for precision and recall probes.
    \item \textbf{Negative Sampling (Precision Probes):} 10,000 instances. Discuss operators like \texttt{C3\_short\_long} (vowel length swaps) and \texttt{C9\_old\_accept} (traps for Model A).
    \item \textbf{Positive Sampling (Recall Probes):} 1,200 instances. Discuss "Oracle-blind" positives (e.g., ศัตรู, กษัตรี) intentionally designed to deduct points from Model B where it fails linguistically.
\end{itemize}

\section{Evaluation and Results}
% [Drafting note: The heavy lifting. Use tables here based on PROGRESS.md.]
\subsection{Experimental Setup}
\begin{itemize}
    \item Define the evaluated models (A, B, C, D\_w2p, D\_ssg) and metrics (Precision, Recall, F1, Coverage).
    \item Mention computational mitigations for Model C (persistent worker pool, G2P caching).
\end{itemize}

\subsection{Gold Recall Analysis}
\begin{itemize}
    \item Present the Stanza-level gold recall table.
    \item Highlight that D\_ssg achieves the highest recall (88.5\%) followed by B (87.5\%), vastly outperforming A (73.4\%).
\end{itemize}

\subsection{Augmentation and Error Analysis}
\begin{itemize}
    \item Present the Augmentation-only overall table.
    \item \textbf{Precision vs. Recall Trade-off:} Explain that Model B acts as a perfect oracle (0 FPs on negatives) but misses oracle-blind positives. Model C catches rhymes but has 3,129 FPs (mostly on \texttt{C3\_short\_long}).
    \item \textbf{The Override Advantage:} Show how Model D\_w2p recovers 95.7\% of the oracle-blind spots (e.g., reading เพชร $\rightarrow$ เพ็ด) that Model B misses.
\end{itemize}

\section{Applications and Future Work}
% [Drafting note: Connect back to the overarching narrative.]
\subsection{Educational Impact}
\begin{itemize}
    \item Discuss the interactive Hugging Face web interface.
    \item Explain how deterministic, explainable feedback aids teachers and students in learning Klon 4/8 structures without LLM hallucinations.
\end{itemize}

\subsection{Foundation for Generative AI}
\begin{itemize}
    \item Discuss how this verifier is currently implemented as the Reward Model environment for Reinforcement Learning (RL) to train a small LLM to write Klon-Paed.
\end{itemize}

\section{Conclusion}
\begin{itemize}
    \item Summarize: The hybrid approach effectively marries the reliability of rule-based logic with the edge-case coverage of G2P overrides.
    \item Final impact statement: The successful integration of Model B into PyThaiNLP 5.3.5 sets a new standard for Thai computational linguistics, paving the way for both educational tools and advanced LLM pipelines.
\end{itemize}

\section*{Acknowledgment}
\begin{itemize}
    \item Acknowledge vajirayana.org for the literary datasets.
    \item Acknowledge the PyThaiNLP maintainers for reviewing and merging the PR.
    \item Acknowledge the developers of the \texttt{tltk} and KlonSuphap-LM baselines.
\end{itemize}

\begin{thebibliography}{00}
% [Drafting note: You will need to add proper citations here later.]
\bibitem{pythainlp} PyThaiNLP Contributors, "PyThaiNLP: Thai Natural Language Processing in Python," [Online]. Available: https://github.com/PyThaiNLP/pythainlp
\bibitem{klonsuphap} Kongfha, "KlonSuphap-LM," [Online]. Available: [Insert Link]
\bibitem{isai_nlp} [Add other relevant Thai NLP, G2P, or LLM RLHF papers here]
\end{thebibliography}

\end{document}